% Feature module: index

\ifdefined\NextFeatureIndexLoaded
  \endinput
\fi
\def\NextFeatureIndexLoaded{1}

\RequirePackage[xindy]{imakeidx}
\UseFeature{hyperlinks}

\providecommand{\IndexTitle}{}
\makeatletter
% xindy receives an opaque occurrence ordinal, never the visible page label.
% The custom location class also covers native \index occurrences without a
% formatting encapsulator.  This preserves the original sort/display grammar.
% Keep the style beside the job (no additional installed xindy search path).
\newwrite\NextIndexStyleFile
\immediate\openout\NextIndexStyleFile=\jobname-impe.xdy
\immediate\write\NextIndexStyleFile{(define-attributes ("NextIndexLocation"))}
\immediate\write\NextIndexStyleFile{(define-location-class "impe-occurrence" ("arabic-numbers" :sep "." "arabic-numbers") :var)}
\immediate\write\NextIndexStyleFile{(markup-locref :class "impe-occurrence" :open "\@backslashchar NextIndexLocation\@charlb" :close "\@charrb")}
\immediate\write\NextIndexStyleFile{(markup-locref :attr "NextIndexLocation" :open "\@backslashchar NextIndexLocation\@charlb" :close "\@charrb")}
\immediate\closeout\NextIndexStyleFile
\DeclareRobustCommand{\NextIndexLocation}[1]{%
  \hyperref[impe-index-#1]{\pageref*{impe-index-#1}}%
}
% Write an occurrence through the shared Hyperlinks destination allocator.
% Grouping prevents the synthetic index label from replacing the most recent
% section/equation/figure reference context for a later user \label.
\newcommand{\NextIndexWriteOccurrenceTo}[2]{%
  \begingroup
    \NextHyperlinksCreateOccurrenceTarget\NextIndexTarget
    \edef\NextIndexOrdinal{0.\the\NextHyperDestinationSequence}%
    \edef\@currentHref{\NextIndexTarget}%
    \label{impe-index-\NextIndexOrdinal}%
    \imki@wrindexentry{#1}{#2|NextIndexLocation}{\NextIndexOrdinal}%
  \endgroup
}
\newcommand{\NextIndexWriteOccurrence}[1]{%
  \NextIndexWriteOccurrenceTo{\jobname}{#1}%
}
\ifx\IndexTitle\@empty
  \makeindex[intoc,options=-M \jobname-impe]
\else
  \makeindex[title={\IndexTitle}, intoc,options=-M \jobname-impe]
\fi
\makeatother

% Native \index: keep the imakeidx public syntax and index grammar.  For the
% default index, an ordinary entry (including sort@display and nested ! keys)
% gets an occurrence-specific destination, just like \Term.
% Explicit |... encapsulations, index ranges/cross-references, and additional
% named indexes retain their native handling until the Index integration (#20)
% establishes an unambiguous formatting/processing contract.
\makeatletter
\newcommand{\NextIndexNativeWrite}[1]{%
  \@bsphack
  \NextIndexWriteOccurrenceTo{\jobname}{#1}%
  \@esphack
}
\makeatother
\ExplSyntaxOn
\cs_new_protected:Npn \__next_index_native_dispatch:n #1
  {
    \tl_if_in:nnTF {#1} {|}
      { \NextIndexNativeIndex{#1} }
      { \NextIndexNativeWrite{#1} }
  }
\cs_new_protected:Npn \NextIndexNativeDispatch #1
  { \__next_index_native_dispatch:n {#1} }
\ExplSyntaxOff
\AtBeginDocument{%
  \let\NextIndexNativeIndex\index
  \RenewDocumentCommand{\index}{o m}{%
    \IfNoValueTF{#1}
      {\NextIndexNativeDispatch{#2}}
      {\NextIndexNativeIndex[#1]{#2}}%
  }%
}

\indexsetup{
  othercode = {
    \setlength{\parindent}{0pt}
    \setlength{\parskip}{0pt}
    \linespread{1.10}\selectfont
  }
}

\DeclareRobustCommand{\NextIndexCjkDescription}[1]{（#1）}
\DeclareRobustCommand{\NextIndexWesternDescription}[1]{%
  \nobreakspace(#1)%
}
\DeclareRobustCommand{\NextIndexBareDescription}[1]{%
  \nobreakspace#1%
}

\ExplSyntaxOn

\prop_new:N \g__next_index_seen_prop
\tl_new:N \l__next_index_parentheses_tl
\tl_new:N \l__next_index_description_tl
\tl_new:N \l__next_index_display_tl
\tl_new:N \l__next_index_sort_tl
\tl_new:N \l__next_index_key_tl

\keys_define:nn { next/index/term }
  {
    sort .tl_set:N = \l__next_index_sort_tl,
    key .tl_set:N = \l__next_index_key_tl,
    parentheses .choice:,
    parentheses / auto .code:n =
      { \tl_set:Nn \l__next_index_parentheses_tl { auto } },
    parentheses / cjk .code:n =
      { \tl_set:Nn \l__next_index_parentheses_tl { cjk } },
    parentheses / western .code:n =
      { \tl_set:Nn \l__next_index_parentheses_tl { western } },
    parentheses / none .code:n =
      { \tl_set:Nn \l__next_index_parentheses_tl { none } },
    parentheses .value_required:n = true,
    parens .meta:n = { parentheses = #1 },
  }

\cs_new_protected:Npn \__next_index_resolve_parentheses:
  {
    \str_if_eq:VnT \l__next_index_parentheses_tl { auto }
      {
        \tl_set:Nn \l__next_index_parentheses_tl { western }
        \cs_if_exist:NT \NextSystemUiTag
          {
            \str_if_eq:eeT { \NextSystemUiTag } { zh }
              { \tl_set:Nn \l__next_index_parentheses_tl { cjk } }
          }
      }
  }

\cs_new_protected:Npn \__next_index_set_description:n #1
  {
    \tl_clear:N \l__next_index_description_tl
    \tl_if_blank:nF {#1}
      {
        \str_case:Vn \l__next_index_parentheses_tl
          {
            { cjk }
              {
                \tl_set:Nn \l__next_index_description_tl
                  { \NextIndexCjkDescription{#1} }
              }
            { western }
              {
                \tl_set:Nn \l__next_index_description_tl
                  { \NextIndexWesternDescription{#1} }
              }
            { none }
              {
                \tl_set:Nn \l__next_index_description_tl
                  { \NextIndexBareDescription{#1} }
              }
          }
      }
  }

\cs_new_protected:Npn \__next_index_emit_term:nnnn #1#2#3#4
  {
    \textbf{#3}#4\unskip\nobreak\hspace{0pt}
    \prop_if_in:NnF \g__next_index_seen_prop {#1}
      {
        \prop_gput:Nnn \g__next_index_seen_prop {#1} { indexed }
        \NextIndexWriteOccurrence{#2@#3#4}
      }
  }
\cs_generate_variant:Nn \__next_index_emit_term:nnnn { VVVV }

\NewDocumentCommand{\Term}{ O{} m o g g }
  {
    \group_begin:
      \tl_set:Nn \l__next_index_parentheses_tl { auto }
      \tl_set:Nn \l__next_index_display_tl {#2}
      \tl_set:Nn \l__next_index_sort_tl {#2}
      \tl_set:Nn \l__next_index_key_tl {#2}
      \tl_clear:N \l__next_index_description_tl
      \IfNoValueTF{#5}
        {
          \IfNoValueTF{#3}
            {
              % Also accept the recently introduced braced optional form.
              \IfNoValueF{#4}
                { \tl_set:Nn \l__next_index_description_tl {#4} }
            }
            { \tl_set:Nn \l__next_index_description_tl {#3} }
        }
        {
          % Legacy form: \Term{key}{display}{description}
          \tl_set:Nn \l__next_index_display_tl {#4}
          \tl_set:Nn \l__next_index_description_tl {#5}
        }
      \keys_set:nn { next/index/term } {#1}
      \__next_index_resolve_parentheses:
      \exp_args:NV \__next_index_set_description:n
        \l__next_index_description_tl
      \__next_index_emit_term:VVVV
        \l__next_index_key_tl
        \l__next_index_sort_tl
        \l__next_index_display_tl
        \l__next_index_description_tl
    \group_end:
  }

\ExplSyntaxOff
